\documentclass[9pt,a4paper]{extarticle}
\usepackage{parskip}

% ======================
% Packages
% ======================
\usepackage[utf8]{inputenc}
\usepackage[T1]{fontenc}
\usepackage[brazil]{babel}
\usepackage{amsmath, amssymb}
\usepackage{amsthm}
\usepackage{geometry}
\usepackage{listings}
\usepackage{xcolor}
\usepackage{hyperref}
\usepackage{booktabs}
\usepackage{array}
\usepackage{tabularx}
\usepackage{enumitem}
\usepackage{microtype}
\usepackage{csquotes}
\usepackage{natbib}

\definecolor{important}{RGB}{255, 100, 100}
\definecolor{notice}{RGB}{0, 102, 204}

\newcommand{\impt}[1]{\textcolor{important}{#1}}
\newcommand{\ntc}[1]{\textcolor{notice}{#1}}

\setlist{nosep}
\geometry{margin=2cm}

% ======================
% Listings setup (Python)
% ======================
\definecolor{codebg}{rgb}{0.95,0.95,0.95}
\lstset{
    language=Python,
    backgroundcolor=\color{codebg},
    basicstyle=\ttfamily\footnotesize,
    keywordstyle=\color{blue},
    commentstyle=\color{gray},
    stringstyle=\color{red!70!black},
    showstringspaces=false,
    frame=single,
    breaklines=true,
    inputencoding=utf8
}

% ======================
% Document
% ======================
\begin{document}

\title{Rage Against the Indutivism\footnote{Material baseado em \citep{chalmers1993ciencia}}}
\author{BCC502 -- Metodologia Científica para Ciência da Computação}
\date{\today}
\maketitle

\tableofcontents

\section{Recapitulando o Indutivismo}

\section{O Problema da Indução}

\section{Recuo para Probabilidade}

\section{Respostas ao Problema da Indução}

\section{A Dependência que a Observação tem da Teoria}

\subsection{Explicação Popular da Observação}

\subsection{}

\bibliographystyle{apalike}   
\bibliography{references}

\end{document}

